\documentclass{article}
\usepackage[preprint]{tmlr}
\usepackage{graphicx}
\usepackage{booktabs}
\usepackage{amsmath}
\usepackage{amssymb}
\usepackage{longtable}
\usepackage{array}
\usepackage{calc}
\providecommand{\real}[1]{#1}
\usepackage[T1]{fontenc}
\usepackage{microtype}
\title{Same recipe, different model: cross-system reproducibility of byte-level language-model pretraining holds at 10M but fails at 50M parameters}
\author{Tadanobu Chuyo Kamijo\\Department of Systems Physiology, Graduate School of Medicine\\University of the Ryukyus, Okinawa, Japan\\\texttt{tadanobu@cs.u-ryukyu.ac.jp} \quad ORCID: 0000-0001-6587-7231}
\def\month{07}\def\year{2026}
\begin{document}
\maketitle
\begin{abstract}
Neural-network training is expected to be reproducible: the same recipe---identical seed, data, and code---should yield the same model. In practice, results differ across hardware, but the phenomenon is poorly quantified at scale because repeating large-model pretraining many times is prohibitively expensive. We pretrained byte-level Transformer language models (10M--124M parameters) on a 300\,MB permissively licensed (CC-BY/CC0) neuroscience open-access corpus, on a zero-cost heterogeneous fleet spanning an NVIDIA CUDA GPU, an Apple-Silicon Metal (MPS) backend, and CPUs across two machines, sweeping seeds, model sizes, and numerical precisions for 4{,}000 steps. Hypotheses, thresholds, and an anti-rescue ledger were cryptographically sealed before any run; cross-backend divergence was measured as held-out next-token prediction disagreement, with same-seed replicates as determinism controls. Within a machine and backend, CPU and CUDA reproduced bit-identically while MPS was self-non-reproducible (0.07\%). Across systems, divergence emerged with scale and then saturated: cross-backend prediction disagreement rose from 0.1\% at 10M to 11.1\% at 30M and plateaued near 13--14\% at 50M and 75M (three seeds each), with divergent models sharing near-identical validation loss ($|\Delta|=0.012$)---an ``equally-good-but-different'' outcome that aggregate metrics hide. In single-seed probes, lower precision brought divergence onset earlier (fp16 by $\le$ step 50, bf16 by step 400). Our original two-size persistence hypothesis was refuted at 10M, but a pre-registered denser follow-up (30M and 75M added after peer review) confirmed a monotonic, then saturating, scale-dependence. The divergent points confound backend with machine and build, so we attribute them to a cross-system difference rather than the numerical path alone. Small-model, many-run studies on open corpora make this failure mode measurable and reproducible.
\end{abstract}

\section{Introduction}
Reproducibility is a load-bearing assumption of computational science: an independent
re-execution of the same procedure should reproduce the same result. For neural-network
training this assumption is routinely invoked when a ``model'' is treated as a deterministic
function of its recipe --- the training data, the random seed, the code, and the
hyperparameters. In practice, however, the mapping from recipe to trained weights also
passes through a hardware- and library-specific numerical path, and that path is not
bit-neutral. Floating-point addition is not associative, reduction orders differ between
CPU, GPU, and accelerator kernels, and several high-performance kernels are non-deterministic
by default \citep{goldberg1991,pytorch_repro,nvidia_determinism}. The practical
question is not whether such differences exist but whether they matter --- whether they stay
in the numerical noise floor or grow into behaviorally different models.

This question is under-measured precisely where it matters most. Characterizing run-to-run
and cross-hardware variability requires training the same configuration many times, but at
the scale of contemporary language models a single pretraining run can cost thousands of
GPU-hours, so repetition is rare and variability is usually reported, if at all, as a single
seed's loss curve \citep{bouthillier2021,kaplan2020}. The measurement window that \emph{is}
affordable --- small models trained many times, across heterogeneous hardware --- has been
largely unused. Small models are not merely cheap stand-ins; because the same numerical
non-associativity operates at every scale, they let us ask a sharper question: \emph{as a model
grows, does cross-hardware divergence stay negligible, or does it emerge?}

We approach this with three deliberate design choices. First, the study is \textbf{pre-registered}:
the hypotheses, decision thresholds, feasibility caveats, and an explicit anti-rescue ledger
were written and cryptographically sealed before any full run, so that a surprising result
cannot be retrofitted into a confirmed prediction. Second, it is \textbf{zero-cost}: all training
runs on a small heterogeneous fleet already on hand --- an NVIDIA consumer GPU (CUDA), an
Apple-Silicon machine (Metal/MPS), and CPUs on two different microarchitectures --- so the
heterogeneity that is usually a nuisance becomes the independent variable. Third, it is
demonstrated on an \textbf{open, permissively licensed domain corpus}: roughly 300 MB of
CC-BY/CC0 neuroscience open-access full text, so that both the corpus and the resulting
artifacts can be released for exact re-derivation.

Our contribution is threefold. (i) With pre-registered thresholds and three-seed error bars,
we measure cross-system divergence in byte-level language-model pretraining across \textbf{four model
sizes} (10M--75M) and find it \textbf{strongly size-dependent}: negligible at 10M (0.1\% prediction
disagreement), it \textbf{emerges sharply between 10M and 30M} (to 11.1\%) and then \textbf{saturates}
near 11--14\% through 75M. The monotone scale-dependence was itself pre-registered (in a second
sealed experiment prompted by review) and supported. (ii) We show this divergence is \textbf{hidden by aggregate
metrics} --- models disagreeing on 13\% of held-out next-token predictions can have
near-identical validation loss --- and that, in single-seed probes, it is \textbf{worsened by low
precision}, with bf16 and fp16 bringing onset hundreds to thousands of steps earlier.
(iii) We provide a \textbf{reproducibility apparatus} --- a sealed pre-registration with an
anti-rescue ledger, a determinism-controlled harness, an openly licensed corpus with recorded
provenance, and per-run environment fingerprints --- that makes the phenomenon measurable on
commodity hardware and directly re-runnable by others.


\section{Methods}
\subsection{Corpus}\label{corpus}

We built a domain corpus from the NCBI PubMed Central (PMC) Open Access subset via the
E-utilities API (\texttt{esearch}/\texttt{efetch}), restricted to neuroscience full text. Because the
\texttt{open\ access{[}filter{]}} still admits non-commercial and no-derivatives licenses, every article
was gated on its JATS \texttt{\textless{}license\textgreater{}} element and \textbf{only permissive licenses were kept --- CC0,
CC-BY, and CC-BY-SA} (CC-BY-SA was permitted by the gate but none appeared in the final
corpus) --- with all rejections logged for auditability. Article bodies were
extracted from the JATS XML (tables, figures, reference lists, and mathematics removed),
concatenated in ascending PMCID order with a document separator, and stored as a raw
byte stream. The final corpus is 300.12 MB (6,961 documents: 6,933 CC-BY, 28 CC0; zero
NC/ND contamination) with recorded SHA-256, split 99.5/0.5 into train and validation. The
collector is stdlib-only and rate-limited to NCBI's guidance ($\le$3 requests/s). Corpus
assembly is deterministic, so the byte stream is reproducible from the document set.

\subsection{Model and training harness}\label{model-and-training-harness}

Models are byte-level (vocabulary 256) decoder-only Transformers. Three size presets were
used; the preset names (10M/50M/124M) are nominal size-class labels, and the actual parameter
counts are given below (all hyperparameters are stated for exact reproducibility):

{\def\LTcaptype{none} % do not increment counter
\begin{longtable}[]{@{}
  >{\raggedright\arraybackslash}p{(\linewidth - 10\tabcolsep) * \real{0.1304}}
  >{\raggedleft\arraybackslash}p{(\linewidth - 10\tabcolsep) * \real{0.1739}}
  >{\raggedleft\arraybackslash}p{(\linewidth - 10\tabcolsep) * \real{0.1739}}
  >{\raggedleft\arraybackslash}p{(\linewidth - 10\tabcolsep) * \real{0.1739}}
  >{\raggedleft\arraybackslash}p{(\linewidth - 10\tabcolsep) * \real{0.1739}}
  >{\raggedleft\arraybackslash}p{(\linewidth - 10\tabcolsep) * \real{0.1739}}@{}}
\toprule\noalign{}
\begin{minipage}[b]{\linewidth}\raggedright
preset (nominal)
\end{minipage} & \begin{minipage}[b]{\linewidth}\raggedleft
d\_model
\end{minipage} & \begin{minipage}[b]{\linewidth}\raggedleft
layers
\end{minipage} & \begin{minipage}[b]{\linewidth}\raggedleft
heads
\end{minipage} & \begin{minipage}[b]{\linewidth}\raggedleft
block (context)
\end{minipage} & \begin{minipage}[b]{\linewidth}\raggedleft
actual params
\end{minipage} \\
\midrule\noalign{}
\endhead
\bottomrule\noalign{}
\endlastfoot
10M & 384 & 6 & 6 & 256 & 10.9M \\
30M & 512 & 10 & 8 & 256 & 31.9M \\
50M & 640 & 10 & 10 & 256 & 49.7M \\
75M & 704 & 12 & 11 & 256 & 72.0M \\
124M & 768 & 16 & 12 & 512 & 114.2M \\
\end{longtable}
}

(The ``124M'' label follows the GPT-2 naming convention but is not an exact GPT-2 configuration;
our preset has 114.2M parameters. All comparisons use the actual models above.) Training used
AdamW (lr 3e-4), batch 16, gradient
clipping 1.0, and a fixed data-sampling order independent of the model seed, for 4,000 steps.
The harness records, at eight milestone steps (50, 100, 200, 400, 800, 1,600, 3,200, 4,000),
the validation loss, the argmax next-token predictions on a fixed 4,096-token held-out probe,
a SHA-256 fingerprint of all weights, and per-parameter statistics. Every run's output JSON
records the exact PyTorch version, so any accidental interpreter mismatch is self-evident.

Determinism was requested via \texttt{torch.use\_deterministic\_algorithms}, \texttt{cudnn.deterministic},
and \texttt{CUBLAS\_WORKSPACE\_CONFIG} (set before the framework import). We note that CUDA
memory-efficient attention warns that it is non-deterministic; we therefore treated
within-machine, within-backend bit-identity as an empirical control to be verified rather
than assumed (see Results).

\subsection{Heterogeneous fleet}\label{heterogeneous-fleet}

All computation used hardware already on hand at no marginal cost: an NVIDIA GeForce RTX 3060
(CUDA 12.6, driver 560.94; Windows) providing the CUDA backend and an AMD-based CPU; and an
Apple-Silicon Mac mini (macOS) providing the Metal Performance Shaders (MPS) backend and an
ARM CPU. Both machines ran the same PyTorch release (2.12.1) but necessarily
different platform builds --- \texttt{2.12.1} (Apple/MPS) versus \texttt{2.12.1+cu126} (CUDA) --- a difference
we treat as part of the cross-system confound rather than claim away as ``identical.'' Feasibility was hardware-bounded and reported honestly: the 124M model fit only
on CUDA (MPS exhausted 16 GB unified memory; CPU was too slow for 4,000 steps), so the
cross-backend comparison spans 10M and 50M, with 124M contributing within-backend controls.

\subsection{Pre-registration and anti-rescue}\label{pre-registration-and-anti-rescue}

Before any full run, we sealed a pre-registration document (hypotheses H1--H4, decision
bands, feasibility caveats, and a seven-item anti-rescue ledger) by committing it with its
own SHA-256 recorded in the commit message, establishing a time-gate. The ledger explicitly
forbade, among other moves, shifting thresholds to fit results, dropping disagreeing seeds,
extending training until divergence appeared, and re-scoping hypotheses post hoc. All
verdicts below are reported against the sealed thresholds without modification.

\subsection{Divergence metric and hypotheses}\label{divergence-metric-and-hypotheses}

Cross-backend divergence between two runs is the percentage of the 4,096 fixed-probe
positions at which their argmax next-token predictions differ, at a given step; the
final-step value is the primary outcome and is averaged over seeds \{0, 1, 2\}. The
pre-registered hypotheses were: \textbf{H1 (persistence)} --- CUDA$\leftrightarrow$CPU final disagreement $\ge$10\% at
both 10M and 50M (three-band verdict: $\ge$10\% supported / 5--10\% diminished-but-survives / \textless5\%
washed out); \textbf{H2 (onset)} --- the step at which divergence exceeds 1\% moves earlier with
model size; \textbf{H3 (precision)} --- bf16/fp16 reach onset earlier than fp32; and \textbf{H4 (loss
masking)} --- models differing by $\ge$10\% in predictions differ by \textless0.05 in validation loss.
Because 50M has no feasible CPU run, its cross-backend pair is CUDA$\leftrightarrow$MPS, a substitution
recorded at seal time. The determinism control required same-machine, same-backend,
same-seed replicates to be bit-identical (with MPS self-non-reproducibility a documented
expected exception).


\section{Results}
We trained 27 models to 4,000 steps each (Mac: 11; NVIDIA host: 16), all recording the
expected PyTorch build and full eight-point milestone curves. Results are reported against
the sealed thresholds. Unless stated otherwise, cross-backend disagreement is the three-seed
mean; onset and precision traces (H2, H3) are \textbf{single-seed (seed 0)} and are labelled as
such throughout.

\subsection{Determinism controls hold; MPS is the exception}\label{determinism-controls-hold-mps-is-the-exception}

Same-machine, same-backend, same-seed replicates were \textbf{bit-identical} on both the CPU and
the CUDA backends (weight SHA-256 match; 0.00\% prediction disagreement at step 4,000). We
confirmed this empirically rather than assuming it: despite the CUDA memory-efficient-attention
non-determinism warning, two seed-0 CUDA runs produced identical weight fingerprints, so
deterministic kernels were \textbf{achieved}, not merely requested. The \textbf{Metal (MPS) backend was
self-non-reproducible}: two identical-seed runs disagreed on 0.07\% of predictions with
different weight fingerprints, matching the pre-registered expectation. Cross-backend
differences below are therefore read against a CPU/CUDA baseline that is itself exactly
reproducible.

\subsection{H1: divergence emerges with scale, then saturates (persistence refuted; scale-dependence pre-registered and supported)}\label{h1-divergence-emerges-with-scale-then-saturates-persistence-refuted-scale-dependence-pre-registered-and-supported}

At \textbf{10M parameters}, all six cross-backend/cross-machine pairs agreed almost completely:
CUDA$\leftrightarrow$CPU disagreement was \textbf{0.2\%} (three seeds), and every other pair (CPU$\leftrightarrow$CPU across the two
machines, CUDA$\leftrightarrow$MPS, MPS$\leftrightarrow$CPU) fell between 0.1\% and 0.2\% --- in the pre-registered \textbf{washed-out}
band (\textless5\%). At \textbf{50M}, the feasible cross-backend pair, \textbf{CUDA$\leftrightarrow$MPS}, disagreed on
\textbf{13.2\% $\pm$ 1.0\%} of predictions (mean $\pm$ SD over seeds \{0,1,2\}) --- in the \textbf{supported} band
($\ge$10\%). The original sealed \textbf{H1}, which required $\ge$10\% at \emph{both} sizes, is therefore \textbf{refuted
as stated}: reproducibility is not uniform but size-dependent.

Because two sizes cannot characterize a function of scale, we \textbf{pre-registered a second,
denser experiment} (prompted by peer review; sealed before running, commit 3bff416) adding
\textbf{30M} and \textbf{75M} models and predicting (HE1) that CUDA$\leftrightarrow$MPS disagreement increases with size
across 10M/30M/50M/75M. The four-point CUDA$\leftrightarrow$MPS curve (three seeds each) is
\textbf{0.1\% $\rightarrow$ 11.1\% $\pm$ 2.3\% $\rightarrow$ 13.2\% $\pm$ 1.0\% $\rightarrow$ 13.9\% $\pm$ 0.4\%} (10.9M / 31.9M / 49.7M / 72.0M
parameters; Fig. 2). \textbf{HE1 is supported} (weakly monotonic, 75M $\ge$ 50M), and the curve reveals
a \textbf{shape two points could not}: divergence \textbf{emerges sharply between 10M and 30M}
(0.1\% $\rightarrow$ 11.1\%, crossing the 10\% band) and then \textbf{saturates} near 11--14\% from 30M to 75M
(sub-linear). The per-step traces (Fig. 2, right) show 10M stays flat throughout training while
the larger models cross onset by step 400--800. We note the 30M point is the least stable
($\pm$2.3\%; one seed at 8.0\%), consistent with its position at the transition edge; 50M and 75M are
tight ($\pm$1.0\%, $\pm$0.4\%). This scale-dependence is thus a \textbf{pre-registered, supported} finding, not
a post-hoc narrative --- though the divergent points remain CUDA$\leftrightarrow$MPS pairs that confound backend
with machine and PyTorch build, and the saturation mechanism is unresolved (see Discussion).

\subsection{H2: onset present at 50M, absent at 10M (single-seed)}\label{h2-onset-present-at-50m-absent-at-10m-single-seed}

Tracking seed-0 disagreement across steps, the 10M CUDA$\leftrightarrow$CPU pair never crossed the 1\% onset
threshold (final 0.3\%). The 50M CUDA$\leftrightarrow$MPS pair crossed 1\% by \textbf{step 800} (9.2\%) and first
reached the $\ge$10\% band by \textbf{step 1,600}; thereafter it fluctuated non-monotonically between
about 9\% and 12\% (12.3 $\rightarrow$ 10.8 $\rightarrow$ 11.9\% at steps 1,600/3,200/4,000) rather than climbing
smoothly (Fig. 2). With only one size exhibiting any divergence and a single seed, we cannot
fit an onset-versus-size relationship; the direction (larger diverges, smaller does not) is
merely consistent with the pre-registered H2.

\subsection{H3: lower precision brings divergence earlier (supported, single-seed)}\label{h3-lower-precision-brings-divergence-earlier-supported-single-seed}

Holding the CUDA backend and seed (0) fixed at 50M, we compared fp32 against reduced precision
(Fig. 3). Divergence from the fp32 reference appeared earlier and larger at lower precision:
\textbf{fp32$\leftrightarrow$bf16} crossed the 1\% onset by \textbf{step 400}, and \textbf{fp32$\leftrightarrow$fp16} already disagreed on 6.4\%
of predictions by \textbf{step 50} --- i.e.~its onset is at or before our first milestone ($\le$ step 50),
below the resolution of the grid. The fp16 trace is itself strongly non-monotone
(51.8 $\rightarrow$ 45.1 $\rightarrow$ 25.7 $\rightarrow$ 71.9\% at steps 400/800/1,600/4,000); we therefore report the direction
(low precision $\rightarrow$ earlier, larger divergence) as supporting \textbf{H3} but treat the specific
endpoint magnitudes, including the 71.9\% value, as single-seed and not robust. Additional seeds
for the precision axis are pre-registered future work.

\subsection{H4: aggregate loss masks the divergence (supported)}\label{h4-aggregate-loss-masks-the-divergence-supported}

For the 50M CUDA$\leftrightarrow$MPS pair that disagreed on 13.2\% of predictions, the corresponding
validation-loss difference was only \textbf{\textbar $\Delta$\textbar{} = 0.012} --- below the pre-registered 0.05 masking
threshold. The two backends produced models that are \textbf{equally good by validation loss yet
behaviorally different}, supporting \textbf{H4}: a scalar aggregate metric conceals the
cross-system divergence that a prediction-level comparison reveals. This corroborates, on a
CPU/CUDA/Metal backend set and an open corpus, the same-accuracy/different-prediction effect
reported for GPU architectures by prior verifiable-training work \citep{srivastava2024verifiable}; we
do not claim it as novel (see Discussion).

\subsection{Summary}\label{summary}

Within a machine and backend, CPU and CUDA training is exactly reproducible and MPS is not;
across systems, models are reproducible at 10M but diverge at 50M despite matched loss; and,
in single-seed probes, reduced precision moves divergence earlier. Of the four pre-registered
hypotheses, H3 and H4 were supported, H2 was directionally consistent, and H1 was refuted. We
report the refutation as a refutation; the scale-dependence it exposes is developed in the
Discussion as a hypothesis for future, adequately-powered work, not as an established result.


\section{Discussion}
Our result is a pre-registered, four-point scaling curve: on the same 300 MB neuroscience
corpus, the same training recipe produces essentially one model at 10M parameters across CPU,
CUDA, and Metal (0.1\% prediction disagreement) but different models from 30M upward
(11--14\% disagreement) that share the same validation loss. Cross-system divergence \textbf{emerges
sharply between 10M and 30M and then saturates} near 11--14\% through 75M. Two design choices
make this credible. Because the study was pre-registered and sealed, we could not convert a
surprising outcome into a claimed prediction: our original persistence hypothesis (H1) was
refuted at 10M, and the anti-rescue ledger prohibited moving thresholds or excising results;
the scale-dependence itself was then tested under a \textbf{separate pre-registration} (HE1),
prompted by review, and supported. And because within-machine CPU/CUDA training was \emph{verified}
bit-identical, the divergence is not generic run-to-run noise.

\textbf{What we do not claim.} We are careful about the boundaries of the claim. The scale-dependence
is now a \textbf{pre-registered, four-point result} (HE1 supported), not a two-point post-hoc
narrative --- but we still do not claim a \emph{law}: four sizes over one order of magnitude establish
emergence-then-saturation on our corpus and system pair, not a universal exponent, and whether
the plateau holds or moves at larger scale is untested. More importantly, every divergent point
is a \textbf{CUDA$\leftrightarrow$MPS} pair that confounds backend, machine, operating system, and PyTorch platform
build (2.12.1 vs 2.12.1+cu126); the determinism controls exclude run-to-run noise but not these
system axes. We therefore attribute the divergence to a \textbf{difference between two compute
systems}, not to the numerical/hardware path \emph{per se} --- isolating the latter requires a
same-machine cross-backend comparison, which CPU throughput made infeasible here and which is
the first item of future work. Finally, the \textbf{saturation mechanism is unresolved}: why
divergence jumps between 10M and 30M and then plateaus near 11--14\% (rather than growing without
bound) is not explained by our data.

\textbf{Regime versus scale.} Within the real neuroscience corpus and non-overfitting regime, the
four-point curve is monotone and well-behaved (0.1\% $\rightarrow$ 11.1\% $\rightarrow$ 13.2\% $\rightarrow$ 13.9\% at 11M/32M/50M/72M).
The one point that sits \emph{off} this curve is an early overfitting toy pilot
(0.8M-parameter character-level Shakespeare) that diverged by 17--19\% --- larger than our 11M
real-corpus model. We read this not as non-monotonicity of scale but as a \textbf{regime} effect:
in a memorization regime the sharp loss surface amplifies tiny numerical differences, whereas
our generalization-regime models follow the scale curve above. We find this mechanism plausible
but offer \textbf{no direct evidence} for it (no curvature, gradient-noise, or overfitting-gap
measurement), so we do not assert it; the toy result is reported only as a distinct,
corroborating regime. Instrumenting loss-surface sharpness across matched regimes --- and
explaining the 10M$\rightarrow$30M jump and the subsequent plateau --- is a concrete next study.

\textbf{Practical implication, appropriately bounded.} The combination of H1 and H4 carries the
actionable message. From 30M upward, a developer who validates a pipeline by comparing final
loss across machines would see agreement while the resulting models differ on more than one in
ten predictions. Whether the plateau holds or the masked divergence resumes growing beyond 75M
is untested; but across our four sizes it does not shrink once emerged, and single-seed probes
suggest reduced precision (now standard) brings its onset dramatically earlier. Prudence therefore favors verifying training-pipeline portability at the
level of predictions, not aggregate loss, and pinning the numerical/system path when exact
reproducibility is required.

\textbf{Positioning} (see also Related work). That floating-point non-associativity and
non-deterministic kernels make training hardware-dependent is known {[}goldberg1991; shanmugavelu2024fpna;
zhuang2022; summers2021{]}, and --- most relevant to our H4 --- prior work on verifiable training
already reports that the same recipe on different GPU architectures can yield similar aggregate
accuracy but substantially different predictions \citep{srivastava2024verifiable}. We therefore do
\textbf{not} claim the same-loss/different-predictions phenomenon as novel; our H4 corroborates it on
a different backend set (CPU/CUDA/Metal rather than GPU--GPU) and an open corpus. Our specific,
narrower delta is a \textbf{pre-registered, determinism-controlled, prediction-level} measurement of
how cross-system divergence depends on \textbf{model size} for \textbf{training from scratch} --- an axis
these works do not isolate, and one where inference-side evidence points the \emph{other} way
(token-probability nondeterminism reported size-independent from 270M to 12B \citep{fu2026beyond}).
That contrast is itself a reason the training-scale question is worth posing; we pose it and
answer it with a pre-registered four-point curve. The claim is the disciplined, scale-resolved
measurement, not the existence of hardware-induced divergence.

\textbf{Limitations.} (i) Four sizes over one order of magnitude (10M--75M) characterize the
cross-system curve; 124M fit only on CUDA, and the plateau's behavior beyond 75M is untested.
(ii) Every divergent point is a CUDA$\leftrightarrow$MPS pair, singly confounded with machine and build (above).
(iii) H2 and H3 rest on a single seed. (iv) The emergence/saturation mechanism is unmeasured.
(v) Onset is resolved only at eight milestone steps. None threatens the registered qualitative
result --- reproducible at 10M, divergent and saturating from 30M with matched loss --- but each
bounds how far it can be generalized.

\textbf{Outlook.} Beyond closing these gaps (a second backend at 124M via quantization or gradient
accumulation; a same-machine 50M cross-backend pair; multi-seed precision runs; a loss-surface
probe), the apparatus itself is the durable contribution: a sealed pre-registration, a
determinism-controlled harness, an openly licensed corpus with recorded provenance, and per-run
environment fingerprints, all runnable on commodity hardware at no cost --- an instrument the
community can use to map where pretraining stops being reproducible.


\section{Related work and positioning}
paragraph. Written after a literature search; the honest delta is narrower than an early draft
implied --- the \emph{phenomenon} is known, our contribution is a pre-registered, scale-resolved,
open-corpus measurement of it.)

\textbf{Floating-point non-associativity and numerical reproducibility.} That summation order
changes floating-point results is textbook \citep{goldberg1991}, and its consequences for HPC and
deep-learning reproducibility have been characterized directly \citep{shanmugavelu2024fpna}. Parallel GPU
reductions, kernel non-invariance, and batch-size effects are established implementation-level
sources of run-to-run and cross-device variance \citep{shanmugavelu2024fpna,pytorch_repro,nvidia_determinism}.
We do not re-establish these; we take them as the mechanism and ask a downstream question.

\textbf{Hardware nondeterminism yields behaviorally different models at matched aggregate
performance.} Most directly related to our H4, prior work on verifiable training shows that the
\emph{same} recipe on \emph{different GPU architectures} can produce models with similar aggregate
accuracy yet substantially different predictions, and treats this as an obstacle to be
controlled \citep{srivastava2024verifiable}. Related analyses characterize numerical sources of
nondeterminism in LLM inference at temperature zero \citep{yuan2025llminfer}, and optimization-level
work documents instability and nondeterminism during training \citep{summers2021,zhuang2022}. Our
H4 result --- 13.2\% prediction disagreement at near-identical validation loss --- therefore
\textbf{corroborates}, on a different backend set (CPU/CUDA/Metal rather than GPU--GPU) and an open
corpus, a phenomenon these works already report; we do not claim it as novel.

\textbf{Does it scale, and does training differ from inference?} The axis we probe --- how
divergence depends on model size --- has been examined for \emph{inference}: measuring token-probability
nondeterminism across Gemma-family models from 270M to 12B, prior work finds the magnitude of
nondeterminism largely \emph{independent} of model size \citep{fu2026beyond}. Our four-point curve points
the other way for \emph{training from scratch} --- reproducible at 10M, divergent from 30M and
saturating through 75M --- which we read not as a contradiction but as a distinction:
inference-time nondeterminism and training-from-scratch divergence are different quantities (the
latter compounds over optimization steps), so the fact that one is size-independent and the
other size-dependent is itself informative. Cross-backend reproducibility has also been
attacked at the \emph{deployment} level, with configuration-first frameworks that detect and mitigate
numerical drift when a fixed model is run on CPU/GPU/compiled runtimes \citep{li2025crossbackend};
that setting (a trained model, moved across backends) is distinct from ours (a model \emph{trained}
on each backend). Training-time reproducibility itself has an established taxonomy of software
and hardware nondeterminism sources and mitigations \citep{chen2022trainrepro}, on which we build.

\textbf{Variance and reproducibility methodology in ML.} A methodological literature argues for
treating training variance as a first-class quantity and for reporting it across seeds and
tooling \citep{bouthillier2021,henderson2018,pineau2021}. We adopt that stance and add
pre-registration with an anti-rescue ledger, which is uncommon in this systems-ML setting.

\textbf{Precision.} Reduced-precision training is known to be less numerically stable, with fp32
near-deterministic and bf16/fp16 more variable \citep{shanmugavelu2024fpna,yuan2025llminfer}; our single-seed H3
probe is consistent in direction (lower precision $\rightarrow$ earlier onset) and is offered as
corroboration rather than a controlled precision study.

\textbf{Our delta (one sentence per axis).} Relative to the above: (i) we measure how cross-system
divergence in \textbf{training from scratch} depends on \textbf{model size} --- a question left open by
inference-side results that find nondeterminism size-independent \citep{fu2026beyond} --- with a
\textbf{pre-registered four-point curve} (10M--75M) showing that divergence emerges sharply between
10M and 30M and then saturates; (ii) we do so under a \textbf{sealed pre-registration} with an
anti-rescue ledger, so a refuted hypothesis is reported as refuted; (iii) we compare at the
level of \textbf{held-out predictions} with same-seed \textbf{bit-identity controls}, not aggregate loss
alone; and (iv) we release an \textbf{openly licensed corpus} and per-run environment provenance so
the measurement is exactly re-runnable on commodity hardware. The novelty is the disciplined,
scale-resolved \emph{instrument and measurement}, not the existence of hardware-induced divergence.



\begin{figure}[t]\centering
\includegraphics[width=0.86\textwidth]{fig_abstract_schematic.png}
\caption{Graphical abstract. One training recipe (same seed, data, code) on CPU/CUDA/MPS
yields a near-identical model at 10M (0.2\% prediction disagreement) but divergent
``equally-good-but-different'' models at 50M (13.2\%, equal loss).}
\label{fig:abstract}\end{figure}

\begin{figure}[t]\centering
\includegraphics[width=0.98\textwidth]{fig_h1_scale_emergence.png}
\caption{Cross-system (CUDA$\leftrightarrow$MPS) prediction disagreement and model scale.
\emph{Left:} final-step disagreement vs.\ model size (10.9M/31.9M/49.7M/72.0M parameters);
markers are the mean over seeds $\{0,1,2\}$, error bars $\pm1$ SD, dashed line the
pre-registered 10\% band. Divergence emerges sharply between 10M and 30M, then saturates.
\emph{Right:} disagreement vs.\ training step (seed 0) for the four sizes; the 11M model stays
flat throughout training while the larger models cross onset by step 400--800.}
\label{fig:h1}\end{figure}

\begin{figure}[t]\centering
\includegraphics[width=0.7\textwidth]{fig_h3_precision.png}
\caption{Precision and divergence onset at 50M on CUDA (single seed): fp32$\leftrightarrow$bf16
and fp32$\leftrightarrow$fp16, with fp32 CUDA$\leftrightarrow$MPS for reference. Onset
$\le$ step 50 (fp16) and step 400 (bf16).}
\label{fig:h3}\end{figure}


\section*{Data and code availability}
All code, the sealed pre-registration, the analysis, and the figures are available at
\texttt{https://github.com/tckamijo/repro-pretrain} (archived at Zenodo,
DOI: \texttt{10.5281/zenodo.21428182}). The
corpus is not redistributed as raw text; it is exactly reconstructible from the PMC Open Access
subset using the included collector and verifiable against the recorded SHA-256
(\texttt{9d6b168e...5b0a27}, 300{,}120{,}401 bytes). Per-run outputs (validation loss, held-out
predictions, weight fingerprints, PyTorch build) are provided for all 27 runs, and each figure
ships with a provenance sidecar. A claim-to-source verification worksheet is included.

\section*{Competing interests}
The author declares no competing interests.

\section*{Statement on the use of AI tools}
In preparing this work the author used a large language model (Anthropic's Claude) as an
assistive tool: drafting and editing the manuscript; writing and scaffolding the
data-collection, training-harness, analysis, and figure-generation code; assisting with the
literature search; and running an internal adversarial review. All experiments were executed by
the author on the author's own hardware, and every reported quantitative result was produced by
the analysis code and independently verified by the author against the raw per-run outputs. The
author designed the study, sealed the pre-registration, checked all numbers, claims, and
citations, and takes full responsibility for the content and its scientific validity. No data,
results, or citations were fabricated.

\bibliographystyle{tmlr}
\bibliography{refs}
\end{document}
